\documentclass[conference]{IEEEtran}
% Options possibles : conference, journal, 10pt, a4paper, etc.

\usepackage[utf8]{inputenc}
\usepackage[T1]{fontenc}
\usepackage[english]{babel}
\usepackage{graphicx}
\usepackage{amsmath,amssymb}
\usepackage{booktabs}
\usepackage{hyperref}
\usepackage{cite}
\usepackage{siunitx}

% --- Title and authors ---
\title{Reactive Following with Five Tricycle Robots in Webots and ROS~2}

\author{
\IEEEauthorblockN{JOHANSEN Magnus, SANTOS Júlio, COLOMBU Jules, PINHEIRO Nathan}
\IEEEauthorblockA{
Faculty of Engineering of the University of Porto (FEUP)\\
Porto, Portugal\\
\{email1, email2, email3, email4\}@fe.up.pt % TODO
}
}

\begin{document}
\maketitle

\begin{abstract}
\normalfont
This paper presents a purely reactive multi-robots following behaviour
implemented in Webots~R2025a and ROS~2. Five differential-drive robots,
each with its own constant cruise speed and base curvature, are equipped
with an array of eleven frontal distance sensors
($100^{\circ}$ field of view, \SI{4}{m} range). A robot that detects a
neighbour turns towards it and follows it; when the neighbour is lost,
it resumes its own base curve. The controller uses only the current local
sensor readings: no global localisation, no target memory. We describe
the architecture, the controller and its parameters, and we evaluate the
behaviour over several randomised initial conditions.
% [TODO] Add one sentence with the main quantitative result,
% e.g. mean following duration and number of collisions.
\end{abstract}

\section{Introduction}
Simple reactive rules can produce rich collective behaviour:
Braitenberg's vehicles~\cite{braitenberg1984vehicles}, Brooks' layered
reactive control~\cite{brooks1986robust} and Reynolds' flocking
model~\cite{reynolds1987flocks} all show that coupling sensing directly
to motion, without a world model, is enough to generate coordinated
motion. This work follows that spirit in a small experimental setting.

The challenge is the following. Five robots move in a flat arena, each
with a different constant speed and base curvature. When a robot sees a
neighbour in its frontal sensing region, it must turn towards it and
follow it. The control must be strictly reactive (current local readings
only). The goals are to
\begin{enumerate}
    \item implement the controller,
    \item connect sensing and motion through ROS~2, and
    \item study the collective behaviour that emerges (chains, groups
          and failures).
\end{enumerate}

The contributions are:
\begin{itemize}
    \item a reproducible five-robot Webots/ROS~2 setup with seeded
          random initial poses;
    \item a stateless following controller with a documented parameter
          set; and
    \item an experimental analysis of its behaviour.
\end{itemize}

The rest of the paper is organised as follows. Section~\ref{sec:arch}
describes the architecture, Section~\ref{sec:ctrl} the controller,
Section~\ref{sec:exp} the experiments, and
Sections~\ref{sec:disc}--\ref{sec:conc} discuss the results and
conclude.

% ============================================================
\section{System Architecture}

\subsection{Simulation Environment}
% [TODO] Describe the Webots environment, dimensions, seed, and random generation.

\begin{table}[t]
    \centering
    \caption{Simulation environment metrics extracted from the project.}
    \label{tab:environment}
    \begin{tabular}{lll}
        \toprule
        Metric & Value & Source \\
        \midrule
        Webots version & R2025a & World header \\
        Arena dimensions & $6 \times 6$ m & Floor definition \\
        Arena area & \SI{36}{m^2} & Derived \\
        Basic time step & \SI{32}{ms} & World definition \\
        Controller update rate & \SI{31.25}{Hz} & Derived \\
        Robots & 5 & Launch configuration \\
        Spawn zone & $1.4 \times 1.4$ m & YAML configuration \\
        Minimum spawn distance & \SI{0.5}{m} & YAML configuration \\
        Placement attempts / robot & 1000 & YAML configuration \\
        Default random seed & 18 & YAML configuration \\
        \bottomrule
    \end{tabular}
\end{table}

\subsection{Robot Model}
% [TODO] Describe the PROTO, dimensions, sensors, and actuators.

\begin{table}[t]
    \centering
    \caption{Robot model metrics extracted from the PROTO definition.}
    \label{tab:robot}
    \begin{tabular}{lll}
        \toprule
        Metric & Value & Source \\
        \midrule
        Body dimensions & $0.08 \times 0.06 \times 0.03$ m & PROTO geometry \\
        Wheel radius & \SI{0.02}{m} & PROTO geometry \\
        Wheel centre offset & \SI{0.035}{m} & PROTO joints \\
        Distance sensors & 11 & PROTO definition \\
        Sensor field of view & $100^{\circ}$ & Sensor angles \\
        Sensor range & \SI{4}{m} & Lookup table \\
        Motor maximum velocity & \SI{20}{rad/s} & Motor definition \\
        \bottomrule
    \end{tabular}
\end{table}

\subsection{ROS~2 Integration}
% [TODO] Describe the launch file, namespaces, Python plugin, and YAML parameters.

\begin{table}[t]
    \centering
    \caption{ROS~2 integration metrics.}
    \label{tab:ros}
    \begin{tabular}{ll}
        \toprule
        Metric & Value \\
        \midrule
        Package type & \texttt{ament\_python} \\
        Controller instances & 5 \\
        Command topic type & \texttt{geometry\_msgs/msg/Twist} \\
        Sensor topic type & \texttt{std\_msgs/msg/Float32MultiArray} \\
        Command timeout & \SI{1.0}{s} \\
        \bottomrule
    \end{tabular}
\end{table}

\begin{figure}[t]
    \centering
    % \includegraphics[width=\columnwidth]{figures/architecture.pdf}
    \caption{System architecture.}
    \label{fig:architecture}
\end{figure}

% ============================================================
\section{Reactive Controller}
% [TODO] Introduce the stateless controller and its three situations.

\subsection{Cruise Behaviour}
% [TODO] Describe the cruise behaviour and curvature formula.
\begin{equation}
\begin{aligned}
\omega_{L} &= \min\!\bigl(\omega_{\max},\, s_{k}(1 + \kappa_{k})\bigr),\\
\omega_{R} &= \min\!\bigl(\omega_{\max},\, s_{k}(1 - \kappa_{k})\bigr).
\end{aligned}
\label{eq:cruise}
\end{equation}
% [TODO] Explain the meaning of kappa > 0, kappa < 0, and kappa = 0.

\subsection{Detection and Following}
% [TODO] Describe the detection threshold, distance estimation, and lookup-table inversion.
\begin{equation}
s(\rho) =
\begin{cases}
0, & \rho \leq \rho_{\text{stop}},\\[2pt]
\max\!\bigl(\omega_{\max}\, t,\, 0.3\,\omega_{\min}\bigr), & \rho_{\text{stop}} < \rho < \rho_{\text{slow}},\\[2pt]
\omega_{\max}, & \rho \geq \rho_{\text{slow}}.
\end{cases}
\label{eq:speed}
\end{equation}
% [TODO] Describe the steering law (weighted sum of sensor readings).
\begin{equation}
\begin{aligned}
u &= \operatorname{clip}\!\Bigl(\sum_{i:\,d_i > \theta} \frac{d_i}{1024}\sin\phi_i,\; -1,\; 1\Bigr),\\
\Delta &= g\, u\, s(\rho),\\
\omega_{L} &= \operatorname{clip}(s - \Delta,\, 0,\, \omega_{\max}),\\
\omega_{R} &= \operatorname{clip}(s + \Delta,\, 0,\, \omega_{\max}).
\end{aligned}
\label{eq:steering}
\end{equation}

\subsection{Parameters}
% [TODO] Describe the parameters (Tables~\ref{tab:params} and~\ref{tab:cruise}).

\begin{table}[t]
    \centering
    \caption{Controller parameters (default configuration).}
    \label{tab:params}
    \begin{tabular}{lll}
        \toprule
        Symbol & Meaning & Value \\
        \midrule
        $\theta$        & detection threshold (reading) & 50 \\
        $\rho_{\text{stop}}$ & stop distance          & \SI{0.05}{m} \\
        $\rho_{\text{slow}}$ & slow-down distance     & \SI{0.25}{m} \\
        $\omega_{\max},\ \omega_{\min}$ & wheel speed limits & 15, 1~rad/s \\
        $g$             & turning gain                   & 1.5 \\
        $s_k,\ \kappa_k$ & cruise speed / curvature      & Table~\ref{tab:cruise} \\
        \bottomrule
    \end{tabular}
\end{table}

\begin{table}[t]
    \centering
    \caption{Per-robot cruise settings and resulting motion.}
    \label{tab:cruise}
    \begin{tabular}{lcccccc}
        \toprule
        Robot & $s_k$ & $\kappa_k$ & $\omega_L$ & $\omega_R$ & $v$ (m/s) & $\rho_c$ (m) \\
        \midrule
        robot\_0 & 10.0 & $+0.05$ & 10.50 & 9.50 & 0.200 & 0.700 \\
        robot\_1 &  8.0 & $+0.06$ &  8.48 & 7.52 & 0.160 & 0.583 \\
        robot\_2 &  6.0 & $+0.07$ &  6.42 & 5.58 & 0.120 & 0.500 \\
        robot\_3 &  4.0 & $-0.05$ &  3.80 & 4.20 & 0.080 & 0.700 \\
        robot\_4 &  2.0 & $-0.06$ &  1.88 & 2.12 & 0.040 & 1.167 \\
        \bottomrule
    \end{tabular}
\end{table}

% ============================================================
\section{Experiments and Results}

\subsection{Protocol}
% [TODO] Define the three conditions, duration, and metrics.
% Record every additional seed used for evaluation in Table~\ref{tab:seeds}.

The default configuration uses seed~18. This seed determines the
non-overlapping initial poses and, together with the robot index, the
controller noise sequence. It should therefore be reported for every
reproducibility run. A different seed can be selected at launch with
\texttt{random\_seed:=<value>}; all seeds used in the experiments must be
recorded explicitly in Table~\ref{tab:seeds}.

\begin{table}[t]
    \centering
    \caption{Seed protocol used for the experiments.}
    \label{tab:seed-protocol}
    \begin{tabular}{lll}
        \toprule
        Seed & Role & Status \\
        \midrule
        18 & Default and reproducibility baseline & Configured \\
        -- & Additional evaluation seeds & To be recorded \\
        \bottomrule
    \end{tabular}
\end{table}

\begin{table}[t]
    \centering
    \caption{Results per condition (mean $\pm$ std over $N$ runs).}
    \label{tab:results}
    \begin{tabular}{lcc}
        \toprule
        Metric & No following & Weak turn \\
        \midrule
        Following duration (s) & -- & -- \\
        Angular error (deg)    & -- & -- \\
        Target losses          & -- & -- \\
        Collisions             & -- & -- \\
        \bottomrule
    \end{tabular}
\end{table}

\subsection{Quantitative Results}
% [TODO] Add the quantitative summary (Table~\ref{tab:results}).

\subsection{Qualitative Observations}
% [TODO] Add qualitative observations for several seeds.

\begin{table}[t]
    \centering
    \caption{Outcomes to be recorded for each evaluated seed.}
    \label{tab:seeds}
    \begin{tabular}{ll}
        \toprule
        Seed & Observed outcome \\
        \midrule
        18 & -- \\
        -- & -- \\
        -- & -- \\
        \bottomrule
    \end{tabular}
\end{table}

\begin{figure}[t]
    \centering
    % \includegraphics[width=\columnwidth]{figures/trajectories.png}
    \caption{Trajectories drawn by the five robots for the same initial conditions.}
    \label{fig:trajectories}
\end{figure}

% ============================================================
\section{Discussion}
% [TODO] Discuss chains and groups.
% [TODO] Discuss failure modes.
% [TODO] Discuss parameter effects.

% ============================================================
\section{Conclusions and Limitations}
% [TODO] Recall the contribution.
% [TODO] Discuss the main limitations.
% [TODO] Describe natural extensions.

% ============================================================
\section*{Code and Video}
% [TODO] Add the code repository URL and demonstration video link.

% ============================================================
\section*{Author Contributions}
% [TODO] Add one line per author describing their contribution.

% ============================================================
\nocite{*}
\bibliographystyle{IEEEtran}
\bibliography{references}

\end{document}